\documentclass[11pt,a4paper]{ctexart}
\usepackage[margin=2.35cm]{geometry}
\usepackage{amsmath,amssymb}
\usepackage{enumitem}
\usepackage{xcolor}
\usepackage[most]{tcolorbox}
\usepackage{hyperref}

\hypersetup{
  colorlinks=true,
  linkcolor=black,
  urlcolor=blue!55!black,
  pdftitle={AI4Math 工具概览：Matlas、Rethlas、Archon 与 Danus},
  pdfauthor={秋木}
}
\setlength{\parindent}{2em}
\setlength{\parskip}{0.2em}
\setlist[itemize,enumerate]{topsep=0.3em,itemsep=0.15em}

\newtcolorbox{statementbox}[1]{
  enhanced,
  breakable,
  colback=blue!3,
  colframe=blue!35!black,
  boxrule=0.7pt,
  arc=2mm,
  left=2.5mm,
  right=2.5mm,
  top=1.5mm,
  bottom=1.5mm,
  before skip=0.7em,
  after skip=0.7em,
  fonttitle=\bfseries,
  title={#1}
}

\newcommand{\question}[1]{%
  \par\smallskip
  \noindent{\large\bfseries #1}\par
}

\title{AI4Math 工具概览\\[0.35em]\large Matlas、Rethlas、Archon 与 Danus}
\author{秋木}
\date{2026 年 11 月 7 日}

\begin{document}
\maketitle

\section*{Matlas}

\begin{itemize}
  \item \textbf{Matlas}：面向\textbf{数学陈述}的\textbf{语义检索}，为\textbf{非形式推理}提供\textbf{跨领域定理线索}。
\end{itemize}

解析这句话中的几个概念。首先，从这句话来说，Matlas 像是一个搜索引擎。那么，它搜索的结果是什么呢？和搜索引擎作比较，当用户输入：

\begin{statementbox}{输入问题}
对于一个有限自动机，什么条件能够保证两个状态不可区分？
\end{statementbox}

对于普通搜索引擎，可能返回包含“有限自动机”“状态不可区分”等关键词的\textbf{网页或论文}；Matlas 更希望返回类似这样的\textbf{数学陈述}：

\begin{statementbox}{数学陈述示例}
设 DFA 为 $A=(Q,\Sigma,\delta,q_0,F)$，$\delta^{\ast}$ 是 $\delta$ 在单词上的扩展。两个状态 $p,q\in Q$ 等价，当且仅当
\[
p\sim_A q
\quad\Longleftrightarrow\quad
\forall w\in\Sigma^{\ast},\;
\bigl(\delta^{\ast}(p,w)\in F
\iff
\delta^{\ast}(q,w)\in F\bigr).
\]
\end{statementbox}

以及与它语义相关的：

\begin{itemize}
  \item Myhill--Nerode 等价关系的定义；
  \item DFA 最小化定理；
  \item 可区分状态与区分词的有关引理；
  \item 这些定理所在的论文或教材；
  \item 相关陈述使用了哪些前置定义和定理。
\end{itemize}

更直观地来说，用户向 Matlas 输入类似于上例中的内容（即数学陈述），Matlas 返回\textbf{可能有助于回答问题的已有数学结果}。例如，研究 Agent 正在证明命题 $P$。Matlas 可能返回：

\begin{enumerate}
  \item 一个能直接推出 $P$ 的已知定理；
  \item 一个可以用于证明 $P$ 的中间引理；
  \item 一个表明 $P$ 需要增加假设的反例；
  \item 一个来自其他数学领域、但结构相似的定理；
  \item 一个已经研究过 $P$ 的更一般版本的文献。
\end{enumerate}

\bigskip\hrule\bigskip

下面是我个人在学习过程中的一些问题：

\question{那么用 matlas 来搜索，和把相关提示词输入进基模来进行搜索有什么区别吗？}

主要区别是 Matlas 和基模的行为逻辑不同。基模可以看作是一个模型参数，用户的输入被模型参数映射为输出。而 Matlas 的行为逻辑是：用户输入被 Embedding 模型转化为向量，然后另有一个 Matlas 的数据库，其中的数据都有对应的向量表示，最后通过比较输入与数据库中数据向量的余弦相似度来进行排序，输出数据库中的陈述。

\question{那为什么基模不能模仿 Matlas 的行为呢？或者更直接一点，直接调用 Matlas？}

基模可以模仿 Matlas 的行为逻辑，但是本质上，不调用外部工具的基模只是一段模型参数，人类和基模对话，本质上就是输入经过模型参数的映射成为输出。

当然，在配置了合适的 API 之后，也可以让基模调用 Matlas，那么此时回答的质量体现的就主要是 Matlas 的能力了。

现在回到模型训练的本质上来做一个比较：模型训练就是把数据输入经过训练的映射，输出一个模型参数。所以，如果把 Matlas 比喻成一个有目录的图书馆，基模就是一个读过这个图书馆中所有书籍的人。

\question{在实际测试中，基模和 Matlas 的能力比较如何？}

Matlas 的论文中目前没有给出与\textbf{纯基模}的能力测试。Matlas 的论文主要介绍了其数据库和检索方法，并给出使用案例。

\question{那么这样看来，Matlas 这种工具实际上是提升了基模的能力（如果允许基模调用），那么基模公司是否付费使用 Matlas 工具？}

目前并没有证据表明 OpenAI、Google 和 Anthropic 公司正在付费使用 Matlas。但是，如果未来基模公司决定把 Matlas 内置为官方工具，那么就需要与 Matlas 团队签订相关商业文件。

\question{那么，有没有其他类似工具被基模公司购买的案例？}

有类似案例，但更常见的形式不是“收购整个工具”，而是\textbf{内容授权}或\textbf{工具接入}：

\begin{itemize}
  \item OpenAI 与 Financial Times、News Corp 签订授权协议，让 ChatGPT 使用其当前及历史内容；
  \item Google 通过 Reddit Data API 获得实时、结构化的数据访问；
  \item Claude 接入 Wiley Scholar Gateway，以检索经过同行评审的科研内容；
  \item ChatGPT 和 Claude 都可以接入 Wolfram，让基模调用 Wolfram|Alpha 和 Wolfram Language 进行可靠计算与知识查询。
\end{itemize}

这些案例说明，基模公司确实会通过外部数据库或专业工具增强模型能力。不过，大部分协议只公开合作关系，不公开具体费用，因此不能据此确认基模公司向工具方支付了多少费用。Matlas 目前也没有公开类似的商业合作。

\section*{Rethlas}

\begin{itemize}
  \item \textbf{Rethlas}：自然语言数学研究 Agent，借助 Matlas 探索路线并生成候选证明。
\end{itemize}

从这句描述来看，Rethlas 是一个生成证明的工具，而且存在形式是像 Codex 一样的 Agent。它由两个 Agent 组成：生成 Agent 和验证 Agent。

生成 Agent 读取问题、搜索文献、探索证明思路并生成自然语言的证明；验证 Agent 检查证明逻辑、证明中引用的定理，并把检查结果反馈给生成 Agent（如果发现错误）。

但是值得注意的是，Rethlas 仅生成自然语言证明，而不是 Lean Proof。

\question{多 Agent 相比单 Agent 在能力上有什么提升？}

最重要的提升就是\textbf{并行工作}的能力，即多 Agent 可以让整个系统同时执行相互独立的工作。此外，如果把所有工作都集中在一个 Agent 中，会\textbf{占用大量的上下文}。以及在验证时，让 Agent 验证自己的思路时，不容易发现其中的疏漏。

\section*{Archon}

\begin{itemize}
  \item \textbf{Archon}：形式化 Agent，借助 LeanSearch 将非形式论证转成 Lean 4 项目。
\end{itemize}

\section*{Danus}

\begin{itemize}
  \item \textbf{Danus}：以\textbf{事实图}协调多个推理 worker，并让\textbf{共享记忆受验证器约束}。
\end{itemize}

所谓事实图，就是用图的结构保存数学事实，以及不同数学事实之间的依赖关系。对于一个新写入的证明或者事实，需要经过验证器的检验之后，才能在共享记忆（即事实图）中增添新的点和边。

不过值得注意的是，Danus 目前使用的还是 LLM 验证器，而非形式化验证器。

\bigskip
\noindent\small\textbf{博客原文：}\url{https://qiumu0o0.github.io/posts/AcademicTech/2026.11.07AI4math/}

\end{document}
